\chapter{数据清洗：让AI帮你，但别让AI替你决策}

数据清洗常被误解为删除异常值、填补缺失值和统一格式的技术步骤。更准确地说，它是一组把原始记录转化为可用于特定分析目的之数据的决策。早期数据清洗综述已经指出，模式差异、重复、冲突和错误值往往跨越多个数据源，不能靠单一算法一次解决\cite{rahm2000data}。因此，“清洗完成”不是数据获得了绝对正确的状态，而是当前版本满足了明确用途、规则和质量标准。

数据清洗在整个分析流程中所处的位置往往被低估。许多人把它当作一项机械的预处理步骤——运行几个脚本，填上缺失值，删掉离群点，然后进入“真正的”建模工作。这种理解是错误的。

数据清洗是一系列决策的集合。每一个关于是否填充、是否删除、是否保留某条记录的判断，都涉及对“什么是正确的数据”的定义。而这个定义无法从数据本身推导出来，它来自领域知识、业务逻辑和研究目的。AI 可以在你给出定义之后高效地执行清洗操作，但它无法替你做出这个定义——它不知道你的检测器在故障时速度归零是什么含义，不知道某条高速公路在暴雨天气下的流量峰值是真实信号还是噪声。

本章的核心主张是：\textbf{清洗决策在人，实现在AI，验证仍在人}。这三个环节缺少任何一个，清洗工作都不完整。

% ----------------------------------------------------------------
\section{数据清洗的核心原则：决策在人，实现在AI}
% ----------------------------------------------------------------

\subsection{清洗本质上是一系列领域判断}

同一条记录在不同研究问题下可能需要不同处理。例如，传感器记录的零速度对于常规流量预测可能是设备故障，对于事故检测却可能是最重要的信号。由此可见，清洗规则必须与分析目标绑定；项目也不应只保存一份被覆盖后的“干净数据”，而应保留原始数据、转换脚本、规则说明和质量报告，使每项修改都能够追溯和重现。

从技术角度看，数据清洗不外乎几类操作：识别并处理缺失值、检测并处理异常值、统一格式和编码、解决字段之间的逻辑矛盾。这些操作的执行，AI 可以轻松完成。

但在执行之前，必须先回答一系列问题，而这些问题没有通用答案。\textbf{什么算缺失？}传感器字段值为 \texttt{-1} 代表离线，是否与 \texttt{NaN} 等价？公交刷卡记录中某站点无数据，是真的没有乘客，还是读卡设备故障？\textbf{什么算异常？}速度值 200 km/h 在普通城市道路上肯定是错误，但在某些高速公路路段上呢？凌晨 3 点某区域突然涌现大量订单，是数据错误还是夜班工人下班的真实信号？\textbf{缺失值应该删除还是填充？}如果填充，用什么值？如果删除，删掉的记录会不会让剩余数据的统计特征产生系统性偏差？此外，\textbf{处理结果是否可被追溯？}三个月后，你能否重现清洗步骤并解释每个决策的理由？

这些问题的答案取决于数据来源、分析目的、业务规则和领域惯例——它们都是领域知识，不是统计技巧。把这些问题扔给 AI 解决，得到的是 AI 基于训练数据统计规律猜测的答案，而不是适合你具体场景的正确答案。

\subsection{错误的工作流与正确的工作流}

在实际工作中，有一种常见的错误模式：

\begin{quote}
\textbf{错误工作流（以结果驱动）：}先让 AI 自动清洗数据，再检查清洗后的数据是否“看起来合理”。
\end{quote}

这种工作流的风险在于，AI 在执行清洗时会做出大量隐性假设。例如，它可能默认用列均值填充所有缺失值，用统计阈值删除所有“异常”值。如果你不知道这些假设，你检查结果的时候很可能看不出问题——清洗后的数据看起来很整洁，但背后的决策与你的领域认知完全不符。这类问题在后续建模阶段才会浮现，届时已经难以追溯根源。

正确的工作流是：

\begin{enumerate}
  \item \textbf{先写清洗规则：}在开始任何操作之前，用文字明确描述每类问题的处理方式。这些规则不需要是完整的，但必须覆盖数据中已知的主要问题。
  \item \textbf{让AI按规则生成代码：}把规则用 Prompt 描述清楚，让 AI 生成对应的代码。如果规则本身存在歧义，AI 会在代码中暴露出来，这是一个发现规则不完整的机会。
  \item \textbf{先让AI说明计划，再执行：}在代码生成之前，要求 AI 先用文字描述它打算如何处理每类情况。这一步让你在执行前就能审查 AI 的理解是否与你的规则一致。
  \item \textbf{分步执行，每步验证：}一次只处理一类问题，每步完成后检查结果是否符合预期，再进入下一步。
  \item \textbf{记录所有决策：}哪些规则被应用，处理前后各有多少条记录，删除了哪些，原因是什么。这份记录是分析工作的一部分，不是可选项。
\end{enumerate}

\subsection{AI在清洗工作中的正确定位}

明确了决策在人之后，AI 的价值并不因此变小，而是变得更加清晰。在你已经给出明确规则的前提下，AI 可以快速生成规则对应的代码，包括边界条件处理和日志记录；可以生成清洗前后的对比统计，帮助你验证规则执行的效果；可以识别你尚未意识到的数据质量问题（例如格式不一致、隐藏的重复记录）；可以将一套规则应用于多个数据文件或多个时间段；还可以生成清洗步骤的文档说明，作为分析报告的组成部分。

AI 做不到的是：判断你数据中速度为零是否是有意义的状态，判断某次流量峰值是真实事件还是传感器噪声，判断缺失数据是否应该保留标记而非删除。这些判断必须由你来做，因为它们依赖你对数据来源和业务场景的理解。

% ----------------------------------------------------------------
\section{缺失值：理解缺失机制再决定处理方式}
% ----------------------------------------------------------------

\subsection{缺失机制的三种类型}

对缺失值的处理不能只根据缺失比例决定。缺失数据理论区分完全随机缺失、随机缺失和非随机缺失，其意义在于：缺失发生机制决定删除或填补是否会引入偏差\cite{little2019statistical}。在实践中，这三类机制通常无法仅靠数据本身被完全证明，仍需要结合设备运行、业务流程和采样制度提出合理假设，并通过敏感性分析考察不同假设对结论的影响。

缺失值的处理方式取决于数据为何缺失。统计学提供了三种缺失机制的分类框架，理解这个框架能帮助你在面对具体场景时做出合理的处理决策。

\textbf{完全随机缺失（MCAR）}是指数据是否缺失与观测值本身和其他变量均无关，缺失完全是随机发生的。在这种情况下，删除缺失记录或用简单均值填充不会引入系统性偏差，只会减少样本量。典型例子是 GPS 设备随机性的信号丢失，与车辆位置、速度和时间均无关联。

\textbf{随机缺失（MAR）}是指数据是否缺失与其他已观测变量有关，但与缺失变量本身的真实值无关。在这种情况下，可以利用其他变量来辅助填充，简单删除可能引入偏差。典型例子是检测器数据在特定路段缺失率更高，但这与该路段的路段类型（隧道、立交桥）有关，而与流量本身无关。

\textbf{非随机缺失（MNAR）}是指数据是否缺失与缺失变量本身的真实值有关。这是最难处理的情况，简单的填充方法都会引入偏差。典型例子是网约车平台在高需求时段出现数据记录异常（服务器压力），这意味着流量越大的时段缺失越严重，如果用均值填充，会系统性地低估高峰期需求。

在实际工作中，MNAR 是最危险的情形，因为它会导致所有依赖填充数据的下游分析产生系统性偏差。交通数据中的 MNAR 并不罕见：检测器在极端天气下更容易故障，意味着缺失率最高时恰好是流量最需要被观测的时候。

\subsection{交通数据中的缺失机制例子}

以下列举几类交通领域中常见的缺失场景，它们的缺失机制各不相同，处理方式也因此不同。

\textbf{感应线圈检测器故障：}检测器在极端高温、道路施工或设备老化后更容易故障。这是 MNAR：流量越大的路段，检测器承受的负荷可能越高，故障率可能越高，缺失恰好集中在高流量时段。处理方式是用相邻检测器数据辅助填充，或明确标记不可用时段，而不是简单删除。

\textbf{公交车GPS信号遮挡：}在隧道、高架桥下方和高楼密集区，GPS 信号会中断。遮挡位置是固定的，与车速和时间无关——这更接近 MAR。处理方式是利用轨迹点之间的时间和已知速度进行插值，或利用公交时刻表修正。

\textbf{公交刷卡记录空白：}某些站点在某些时段无刷卡记录。这有两种解读：一是真的没有乘客（MCAR），二是读卡设备故障或数据上传失败（需要调查）。如果直接按零乘客处理，会低估实际需求。处理方式是结合其他信息源（如视频监控、同站历史记录）判断缺失原因，再决定是填充还是标记为无效。

\textbf{出租车轨迹记录中断：}在某些城市，出租车数据存在“空驶段记录缺失”的问题：司机空车时关闭某些上报功能，导致空驶轨迹缺失。这是 MNAR：缺失与司机行为（空驶状态）直接相关。处理方式是不宜填充，应明确记录该数据局限性并在分析中说明。

\subsection{主要处理策略及适用条件}

了解缺失机制之后，可以对应选择处理策略。

\textbf{直接删除}适用于 MCAR 且缺失比例较低（通常低于 5\%）的情形，删除后不会引入明显偏差。它不适用于 MNAR 或缺失集中在特定时段或区域的情形。

\textbf{均值/中位数填充}适用于 MCAR 且该字段用于描述局部水平、不需要保持时序连续性的情形。需要注意的是，填充后会压缩该字段的方差，可能影响后续建模。

\textbf{前向/后向填充}适用于时序数据中短时段的缺失，且相邻值变化较小。在交通数据中，流量平稳的非高峰时段，2--3 个时间步的检测器缺失可以用前向值填充。它不适用于变化快速的时段（如早高峰启动期）。

\textbf{线性插值/样条插值}适用于时序连续且数值变化平滑的场景。对于连续不超过 3 个时间步的检测器数据缺失，线性插值是合理的选择。它不适用于长段缺失或有突变的场景，否则插值会“创造”出不存在的平滑过渡。

\textbf{基于模型的填充}适用于 MAR 且有其他与缺失变量相关的已观测变量可用作预测依据的情形。例如，利用相邻检测器的流量数据，用回归或 KNN 填充故障检测器的缺失值。代价是需要训练模型，且填充质量依赖特征选择。

\textbf{保留并标记}适用于 MNAR、缺失原因不明，或缺失本身是有意义信息的情形。做法是添加一个二值标记字段（例如 \texttt{speed\_is\_missing}），在保留原始缺失值的同时告诉模型该观测点缺失这一事实。这在树模型和神经网络中尤为有用。

\subsection{用Prompt描述缺失值处理规则}

在向 AI 请求生成缺失值处理代码之前，应将处理规则用自然语言明确写出。以下是一个针对道路检测器数据的完整 Prompt 示例：

\begin{quote}
我有一份道路检测器数据，文件名为 \texttt{detector\_raw.csv}，包含以下字段：

\begin{itemize}
  \item \texttt{timestamp}：时间戳，格式为 \texttt{YYYY-MM-DD HH:MM:SS}，5 分钟粒度
  \item \texttt{detector\_id}：检测器编号（字符串）
  \item \texttt{speed}：平均速度（km/h），浮点数
  \item \texttt{volume}：流量（辆/5 分钟），整数
  \item \texttt{status}：检测器状态，取值为 \texttt{normal}（正常）或 \texttt{fault}（故障）
\end{itemize}

请按照以下规则处理缺失值：

\begin{enumerate}
  \item 当 \texttt{status} 为 \texttt{fault} 时，对应行的 \texttt{speed} 和 \texttt{volume} 无论是否为 NaN，均视为无效数据，不进行填充，而是将其标记为 NaN 并在新列 \texttt{speed\_valid} 和 \texttt{volume\_valid}（布尔型）中记录有效性（False 代表无效）；
  \item 当 \texttt{status} 为 \texttt{normal} 且 \texttt{speed} 为 NaN 时：若该检测器当前连续缺失不超过 3 个时间步，使用线性插值填充；若连续缺失超过 3 步，则整段保持 NaN 并在 \texttt{speed\_valid} 中标记为 False；
  \item 对 \texttt{volume} 字段采用相同的缺失处理逻辑；
  \item 请先输出一段文字，描述你对上述规则的理解和处理步骤，确认后再生成代码；
  \item 代码输出时，在每个处理阶段前添加注释说明该步骤的目的；处理完成后输出一个摘要表：每个检测器的原始缺失率、故障导致的无效率、插值填充比例和最终有效数据比例。
\end{enumerate}

请使用 Python，仅使用 pandas 和 numpy。
\end{quote}

注意这个 Prompt 的几个关键做法：首先，规则是分场景的（故障状态 vs 正常状态），而不是统一处理所有缺失值；其次，要求 AI 先描述理解，再执行，这使你有机会在代码生成前发现理解偏差；最后，要求输出验证摘要，使你可以立即检查规则是否被正确执行。

\subsection{验证填充结果}

填充完成后，必须对结果进行验证。验证不是“看看数据是不是整洁了”，而是检查以下具体问题：填充前后各字段的分布是否发生了显著变化，均值、标准差、分位数是否合理；填充值是否引入了不可能出现的数值，例如某段连续低速时段被插值出了一个高速值；标记为无效的记录比例是否与先验认知一致，例如已知某检测器故障率约为 3\%，如果标记结果显示 30\% 无效，需要核查规则是否写错；以及填充后数据在时序图上是否“看起来合理”，突变、平台期和异常曲线是发现填充问题的直观方式。

一个实用的做法是在填充后要求 AI 生成以下输出：每个检测器的速度时序图（含填充前的原始点），以及填充前后的分布对比直方图。这两类图形能快速暴露大部分填充问题。

% ----------------------------------------------------------------
\section{异常值：区分错误与真实极端事件}
% ----------------------------------------------------------------

\subsection{异常值的两种来源}

交通数据中的异常值来源于两类完全不同的原因，处理方式截然相反。

第一类是\textbf{传感器错误或记录错误}，由设备故障、数据传输错误、记录规范不一致或数据合并时的对齐错误产生。这类异常是“脏数据”，应当被识别并清除或修正。典型例子包括：检测器因电压不稳定，周期性输出速度值 999 km/h；GPS 设备时钟漂移，导致轨迹点的时间戳比实际早 5 分钟；手动录入的事故数据中，“死亡人数”字段误录为 99。

第二类是\textbf{真实极端事件}，数据本身是准确的，但因为事件本身超出正常范围，在统计上表现为异常值。这类“异常”是宝贵的信号，删除它们会让模型对真实世界的极端情况失去感知能力。典型例子包括：暴雨天气下某城区出行需求是晴天的 3 倍，在统计上属于极端值，但这是真实的需求波动；事故造成某路段流量骤降至正常值的 10\%，在 IQR 或 Z-score 分析中会被标记为异常，但它是需要被模型捕捉的关键信号；春节期间城市交通量急剧下降，属于季节性极端事件，不应被视为噪声删除。

这一区分的重要性在于：如果把真实极端事件当作错误删除，模型将无法学习这些情形下的规律，在真实部署时面对极端情况时预测会严重失准。

\subsection{为什么不能直接套用统计阈值}

$3\sigma$ 法则和四分位数范围（IQR）方法是教科书中最常见的异常值检测方法。它们有共同的假设：数据的“正常”状态近似服从某种分布，偏离足够多的就是异常。

在交通数据中，这个假设经常不成立。交通流量具有强烈的多峰分布（早高峰、晚高峰、平峰、夜间低谷），不能用单一的均值和标准差描述。工作日和周末、节假日的分布差异巨大，如果混合计算统计阈值，会把正常的工作日高峰标记为异常，或者把正常的节假日低谷标记为异常。真实的极端事件（暴雨、重大赛事、突发事故）确实在统计上是极端值，但它们不是错误。

以一个具体例子说明：某城市道路检测器在工作日早高峰的平均速度为 35 km/h，标准差为 8 km/h。使用 $3\sigma$ 阈值，速度低于 11 km/h 的记录会被标记为异常。但在一次重大暴雪事件中，该路段连续 6 小时速度低于 10 km/h——这是真实的、有意义的极端气候事件，如果将其删除，模型将永远学不到恶劣天气对交通的影响。

\subsection{异常检测策略的层次}

应对交通数据的异常值，推荐的策略不是单一方法，而是分层递进的检测流程：

\begin{enumerate}
  \item \textbf{领域规则过滤（优先级最高）：}使用绝对不可能出现的值作为过滤条件。这些阈值应由领域知识确定，而非统计分布。例如：速度超过 250 km/h 在任何城市道路上都是记录错误；流量值为负数是不可能的；时间戳早于数据采集系统部署日期的记录不可能真实存在。这类规则可以直接删除记录，且无需进一步验证。
  \item \textbf{统计方法初步筛选（生成候选异常列表）：}在排除领域规则明确的错误之后，使用 IQR 或 Z-score 方法识别统计上的异常点。\textbf{关键区别：}这一步的输出是“候选异常”列表，不是“需要删除”列表。统计方法的作用是把值得进一步审查的记录筛选出来，而不是直接决定删除。
  \item \textbf{上下文验证（区分错误与真实事件）：}对候选异常记录进行上下文核查。具体做法包括：与相邻检测器的同时段数据对比（相邻检测器是否也显示异常？）；与历史同期数据对比（同一天同一时段的历史记录是否也有类似值？）；查阅外部记录（是否有当日的气象记录、事故记录、活动记录可以解释这次异常？）。如果多个信息源都指向同一异常，则更可能是真实极端事件，应保留并标记而非删除。
  \item \textbf{人工抽查（质量控制）：}从统计方法标记的候选异常中，随机抽取一部分进行人工逐条核查。抽查比例取决于数据量和时间资源，但即使 5\%--10\% 的抽查也能帮助校正自动检测规则中的系统性偏差。
\end{enumerate}

\subsection{用Prompt描述异常值检测规则}

以下是一个针对公交 GPS 轨迹数据的异常值检测 Prompt：

\begin{quote}
我有一份公交车 GPS 轨迹数据，文件为 \texttt{bus\_gps.csv}，字段包括：

\begin{itemize}
  \item \texttt{bus\_id}：车辆编号
  \item \texttt{timestamp}：记录时间，30 秒粒度
  \item \texttt{lon}、\texttt{lat}：经纬度（WGS84 坐标系）
  \item \texttt{speed}：GPS 推算速度（km/h）
  \item \texttt{heading}：行驶方向（度，0--360）
\end{itemize}

请按以下分层规则检测异常值，不要直接删除，只进行标记：

\begin{enumerate}
  \item \textbf{领域规则异常（标记为 \texttt{error\_domain}）：}\\
    - \texttt{speed} $> 120$（城市公交不可能超过 120 km/h）\\
    - \texttt{speed} $< 0$（速度为负不可能）\\
    - 坐标点超出该城市的合理地理范围（提示：使用以下边界框 [lon\_min, lon\_max, lat\_min, lat\_max]，我稍后填写）\\
    - 相邻两个记录点之间的位移距离超过 500 米（30 秒内移动 500 米以上意味着速度超过 60 m/s，不可能）

  \item \textbf{统计方法候选异常（标记为 \texttt{suspicious\_stat}）：}\\
    - 按车辆分组，计算每辆车的速度 IQR，将 $Q3 + 3 \times \text{IQR}$ 以上的记录标记为候选异常\\
    - 按时间段（0--6 点、7--9 点、10--16 点、17--19 点、20--24 点）分别计算 IQR，而不是对全天统一计算

  \item \textbf{上下文核查辅助信息（生成，不自动处理）：}\\
    - 对每条 \texttt{suspicious\_stat} 记录，输出该记录前后各 3 条记录的速度，供人工核查参考
\end{enumerate}

输出要求：在原数据基础上添加两个布尔列 \texttt{error\_domain} 和 \texttt{suspicious\_stat}，输出到 \texttt{bus\_gps\_flagged.csv}，同时输出一个统计摘要：每辆车的领域规则异常数量和统计候选异常数量。
\end{quote}

\subsection{处理结果的记录与可追溯性}

数据清洗的可追溯性是专业分析的标志之一。清洗决策的记录应包含以下内容：原始数据的基本统计（行数、各字段的缺失率、数值范围）；每类清洗规则的应用结果（有多少记录被删除、有多少被标记、有多少被填充）；删除和标记的原因（是领域规则还是统计方法触发，具体阈值是什么）；清洗后数据的基本统计（用于与原始数据对比，确认变化在预期范围内）；以及关键决策的说明（哪些候选异常经过人工核查后被保留，保留的理由是什么）。

这份记录既是分析报告的基础材料，也是后续复查和方法复现的依据。在撰写论文或技术报告时，数据清洗的记录可以直接转化为“数据预处理”部分的内容，而不需要从头重建。

% ----------------------------------------------------------------
\section{数据格式、类型与一致性问题}
% ----------------------------------------------------------------

格式和一致性问题是数据清洗中最容易被忽视、却最容易引发后续错误的一类问题。它们不像缺失值那么显眼，也不像异常值那么有讨论价值，但一旦存在，可能导致数据合并出错、时序分析失效或空间分析混乱。

\subsection{时间字段的陷阱}

时间字段在交通数据中几乎无处不在，也几乎无处不存在问题。

\textbf{格式不一致}是最常见的问题：同一份数据集中，时间字段可能以 \texttt{2023-04-15 08:30:00}、\texttt{20230415083000}、\texttt{15/04/2023 8:30} 等多种格式混用，如果不统一，日期解析函数会返回错误或产生 NaT。\textbf{时区混乱}在多来源数据合并时尤为危险：不同来源可能使用不同时区，例如出租车平台的时间戳使用本地时间，而 GPS 设备记录的是 UTC 时间，合并后若不统一时区，早高峰的分析结果会错位数小时。\textbf{夏令时跳变}在使用夏令时的地区会导致某些日期存在 23 小时或 25 小时，时序分析中的等时间步假设在这些日期会失效；中国大陆自 1992 年起不再执行夏令时，但如果分析包含其他国家数据或历史数据，这是需要检查的问题。\textbf{采样频率不一致}在数据合并时同样需要预先处理：不同来源的时间粒度可能不同（检测器数据是 5 分钟，气象数据是小时，事件记录是不规则时间戳），在合并前需要明确的上采样或下采样规则，而不是依赖 AI 自行决定对齐方式。

在 Prompt 中描述时间字段问题时，应明确说明：原始格式是什么、目标格式是什么、时区如何统一，以及不同粒度如何对齐。例如：

\begin{quote}
数据中的 \texttt{timestamp} 字段以字符串形式存储，格式混用两种：部分为 \texttt{YYYY-MM-DD HH:MM:SS}，部分为 \texttt{YYYYMMDDHHMMSS}（不含分隔符）。请将两种格式统一解析为 pandas \texttt{datetime64} 类型，时区设定为 UTC+8（北京时间），不使用 \texttt{localize} 方法（因为原始数据未携带时区信息，只需统一解释为北京时间）。统一格式后，输出一个验证摘要，确认是否存在解析失败的记录。
\end{quote}

\subsection{空间字段的陷阱}

坐标系不一致是中国交通数据中最常见的空间数据质量问题，影响范围非常广。

中国目前在实际使用中存在三种坐标系共存的局面：WGS84（国际标准，GPS 设备原始输出使用该坐标系）、GCJ-02（中国国家测绘局加密标准，通常称“火星坐标”，Google 地图和高德地图使用）、BD-09（百度地图专有坐标系，在 GCJ-02 基础上进一步偏移）。三者之间的偏差在城市区域可达数百米，在高精度分析（如路段级匹配）中会导致严重错误。在数据来源与坐标系的对应关系上，GPS 设备原始轨迹通常为 WGS84；高德 API、腾讯地图返回 GCJ-02；百度地图 API 返回 BD-09；国家基础地理信息的 shapefile 通常为 CGCS2000（与 WGS84 在精度要求下可近似等同）。坐标转换是成熟的数学变换，可以让 AI 直接实现，但需要明确说明来源坐标系和目标坐标系：

\begin{quote}
我的出租车 GPS 轨迹数据原始坐标为 WGS84，需要与高德地图路网数据（GCJ-02）进行地图匹配。请实现 WGS84 到 GCJ-02 的坐标转换函数，使用 \texttt{pyproj} 库或手动实现偏移计算均可，并对数据集中所有记录的 \texttt{lon} 和 \texttt{lat} 字段进行转换，将结果存入新列 \texttt{lon\_gcj} 和 \texttt{lat\_gcj}，保留原始列不删除。
\end{quote}

此外，不同数据库中的路段编号不一致也是常见问题。同一条路段，在交管系统中可能编号为 \texttt{R-1023}，在公路管理数据库中为 \texttt{G104-km35}，在 GIS 数据中为 OpenStreetMap 的 way ID。合并多来源数据时，必须建立明确的编号映射关系，而不是依赖 AI 猜测。

\subsection{类别字段的编码不一致}

类别字段的不一致往往比数值字段更隐蔽，因为字符串的差异不会报错，只会悄悄产生错误的分类结果。同一种道路类型，在不同年份的数据中可能分别标记为“城市快速路”、“快速路”和“urban expressway”；同一类事故原因，在不同记录系统中可能有“超速”、“车速过快”、“行驶速度过快”等不同表述；大小写和空格差异同样不容忽视，\texttt{“Normal”}、\texttt{“normal”}、\texttt{“ normal”} 在字符串比较中是三个不同的值；行政区划的表述也会出现差异，同一个地区可能被记录为“南京市”、“江苏南京”或“Nanjing”。

处理这类问题的正确方式是：先通过统计唯一值找出所有变体，再人工定义合并规则，最后让 AI 实现映射：

\begin{quote}
请统计数据中 \texttt{road\_type} 字段的所有唯一值及其出现次数，输出为一个列表，按出现次数从高到低排序。不要做任何合并处理，我将根据这个列表自己定义合并规则，然后告诉你如何合并。
\end{quote}

\subsection{单位不一致}

单位不一致在数值层面不会报错，但会产生严重错误的分析结果。交通数据中常见的单位问题包括：速度字段可能同时存在 km/h 与 m/s（换算系数 3.6）；流量字段可能混用辆/5 分钟、辆/小时和辆/天；密度字段在含有国际数据时可能出现辆/公里与辆/英里并存的情况；距离字段则可能公里与米并存。

在 Prompt 中描述单位要求时，应明确说明目标单位和换算规则，不要假设 AI 知道数据的来源和原始单位：

\begin{quote}
数据中 \texttt{speed\_mps} 字段的单位为 m/s，\texttt{speed\_kmh} 字段的单位为 km/h。请检查这两个字段是否一致（即 \texttt{speed\_kmh} 是否等于 \texttt{speed\_mps} $\times$ 3.6），对不一致的记录输出其比例和样例，并以 \texttt{speed\_kmh} 为准，统一换算后删除 \texttt{speed\_mps} 列。
\end{quote}

% ----------------------------------------------------------------
\section{用Prompt链组织清洗流程}

交互式数据整理工具的研究表明，将用户可理解的变换操作转化为可重复执行的脚本，能够兼顾探索效率与可追溯性\cite{kandel2011wrangler}。与 AI 协作也应遵循这一原则：自然语言用于讨论规则和生成候选实现，最终被接受的处理必须落到确定性脚本、测试和日志中。否则，同一 Prompt 在不同时间产生略有差异的处理结果，清洗过程便无法审计。
% ----------------------------------------------------------------

\subsection{清洗流程的拆解原则}

数据清洗不应一次性完成所有操作。一次性把所有清洗步骤塞入单个 Prompt，不仅难以验证结果，也难以在出错时定位问题所在。

推荐的做法是将清洗流程拆解为若干独立步骤，每步有明确的输入、输出和验证标准。各步骤之间的依赖关系应是线性的：前一步的输出是后一步的输入。一旦某步的输出不符合预期，可以单独修正该步，而不影响其他步骤。

以道路检测器数据的完整清洗流程为例，可以拆解为以下步骤：

\begin{enumerate}
  \item \textbf{初始审查：}读取数据，输出字段列表、数据量、时间范围、每字段的非空数量和数值范围摘要。\textit{验证：时间范围与预期一致，字段数量与数据字典吻合。}
  \item \textbf{格式统一：}统一时间戳格式，检查并修正坐标系（如适用），统一类别字段编码。\textit{验证：时间戳全部可解析，无 NaT；类别字段唯一值数量与业务分类数量一致。}
  \item \textbf{缺失值处理：}按照预定义规则处理各字段的缺失值。\textit{验证：处理后各字段的缺失率，与规则定义的预期相符；填充前后分布对比无重大偏移。}
  \item \textbf{异常值处理：}按照分层规则检测并标记异常值，人工抽查候选异常后决定是否删除。\textit{验证：被删除或标记的记录比例合理，与先验知识一致。}
  \item \textbf{一致性核查：}检查字段之间的逻辑一致性（例如速度为正但流量为零）、跨文件的键值匹配情况。\textit{验证：逻辑矛盾记录比例极低（或能被解释）。}
  \item \textbf{清洗总结：}输出清洗前后对比摘要，记录每步应用的规则和处理结果。\textit{验证：最终有效数据量与分析需求匹配，清洗记录完整。}
\end{enumerate}

\subsection{数据清洗Prompt的模板结构}

针对以上流程，以下提供一个可复用的 Prompt 模板结构，适用于大多数表格型交通数据的清洗任务：

\begin{quote}
\textbf{角色：}你是一位数据工程师，正在处理交通领域的原始数据集，目标是生成可用于后续建模的高质量数据集。

\textbf{数据说明：}数据来源为[描述来源]，文件为[文件名]，行代表[观测单位]，关键字段包括：

\begin{itemize}
  \item [字段名]：[含义]，[单位]，[已知问题或特殊值的含义]
\end{itemize}

\textbf{当前步骤：}[步骤名称，例如“步骤三：缺失值处理”]

\textbf{处理规则：}
\begin{enumerate}
  \item [规则一]
  \item [规则二]
  \item ...
\end{enumerate}

\textbf{约束：}
\begin{itemize}
  \item 不修改输入文件，将输出写入[输出文件名]
  \item 不删除原始字段，新增字段使用[命名约定]
  \item 代码使用 Python，库限于 pandas、numpy[以及其他允许的库]
\end{itemize}

\textbf{执行要求：}请先用文字描述你对每条规则的理解和实现计划，我确认后再生成代码。

\textbf{输出验证：}代码执行完成后，请输出[指定格式的验证摘要]。
\end{quote}

\subsection{关键做法：先让AI描述计划，再执行}

在每个清洗步骤中，要求 AI 先用文字说明它打算如何处理数据，是一个成本低但收益高的做法。

其价值在于：规则的文字描述和代码实现之间往往存在理解差距。AI 可能把“连续缺失不超过 3 步”理解为“缺失总数不超过 3 个”，两者的实现逻辑完全不同。在 AI 输出计划文字时，你可以发现这类理解偏差并立即纠正，避免执行后发现错误再返工。

这一做法的另一个好处是强迫你自己也把规则描述得更清楚。如果 AI 的计划文字与你的预期不符，可能是 AI 理解有误，也可能是你的规则本身存在歧义。两种情况都值得在执行前解决。

\subsection{验证清洗结果的Prompt}

每步清洗完成后，可以用独立的 Prompt 验证结果，而不是把验证逻辑嵌入清洗代码中。分离验证逻辑有助于保持代码的可读性，也便于在验证发现问题时进行针对性修正。

\begin{quote}
清洗步骤已完成，输出文件为 \texttt{detector\_cleaned\_step3.csv}。请对该文件进行以下验证，并以 Markdown 表格形式输出结果，不生成图表：

\begin{enumerate}
  \item 每字段的缺失率（处理后）；
  \item 与上一步输出文件 \texttt{detector\_cleaned\_step2.csv} 相比，每字段的均值、标准差、最小值和最大值的变化；
  \item 检查 \texttt{speed\_valid} 和 \texttt{volume\_valid} 字段的 True/False 比例，与预期规则（约 3\% 的故障率）相符情况；
  \item 任何你在数据中发现的异常模式，请单独列出。
\end{enumerate}

不要推断任何清洗规则，只做统计描述，判断留给我来做。
\end{quote}

最后一句“判断留给我来做”是关键的约束。它明确告诉 AI 不要自作主张地建议删除或修改某些记录，只需要提供客观的统计信息。

\subsection{保留清洗决策记录}

清洗过程的记录应形成一份简洁的决策日志，记录以下要素：每步应用的规则（可直接复制 Prompt 中的规则部分）；每步处理的记录数量变化（原始行数到处理后行数，或标记数量）；关键的人工判断记录（例如：“经人工抽查，2023-08-12 的流量峰值被确认为台风天气引起的真实事件，保留该记录”）；以及最终数据集的基本描述，供论文或报告引用。

这份记录可以用 AI 辅助生成，但应由你确认其准确性。在研究项目中，这份记录的地位与代码同等重要。

% ----------------------------------------------------------------
\section{数据标准化与编码：为建模做准备}
% ----------------------------------------------------------------

数据清洗的终点是建模的起点。在完成缺失值和异常值处理之后，通常还需要对数值特征进行缩放、对时间特征进行编码、对类别特征进行转换，才能作为模型的输入。这些操作在技术上可以交由 AI 实现，但每个决策背后的判断仍然需要人来做。

本节不深入特征工程的全部内容（第六章专门讨论），只梳理与清洗阶段直接相关的标准化和编码决策。

\subsection{数值特征的缩放}

数值特征缩放的目的是消除不同特征量纲带来的影响，防止大数值特征在模型中“天然地”获得更高的权重。常用的两种方法有不同的适用场景。

\textbf{Min-Max标准化（归一化）}将特征值线性映射到 $[0, 1]$ 区间（或其他指定区间），公式为 $x' = (x - x_{\min}) / (x_{\max} - x_{\min})$。它适用于特征有明确上下界的情形（例如道路等级、占用率等），或模型对输入范围敏感的情形（例如神经网络）。当训练集的 $x_{\min}$ 和 $x_{\max}$ 无法代表未来数据范围时（例如流量数据可能在未来出现历史最大值之外的值），这种方法不适合使用。

\textbf{Z-score标准化}将特征转换为均值为 0、标准差为 1 的分布，公式为 $x' = (x - \mu) / \sigma$。它适用于特征近似正态分布，或不确定特征真实范围时（大多数情况下比 Min-Max 更鲁棒）。需要特别注意的是：均值和标准差应该只在训练集上计算，再应用于测试集，不能用全量数据计算，否则会引入数据泄露。

向 AI 描述缩放需求时，必须明确告知：使用哪种方法、在哪个子集上拟合（训练集还是全量）、是否保留原始列，以及如何保存 scaler 对象供后续使用：

\begin{quote}
请对 \texttt{train.csv} 中的 \texttt{speed}、\texttt{volume} 和 \texttt{occupancy} 三个字段进行 Z-score 标准化。拟合参数（均值和标准差）只能在训练集上计算，不使用验证集或测试集的数据。请将标准化后的值存入新列（原列名加 \texttt{\_scaled} 后缀），保留原始列不删除。同时将拟合好的 scaler 对象保存为 \texttt{scaler.pkl}，供后续对测试集 \texttt{test.csv} 进行变换时使用。
\end{quote}

\subsection{时间特征的周期性编码}

时间特征的处理是交通数据分析中一个典型的“看似简单，实则有坑”的问题。

直觉上，把一天中的小时（0--23）直接作为数值特征输入模型，似乎是最简单的做法。但这个做法有一个根本性的问题：它无法表达时间的周期性。

考虑以下场景：凌晨 0 点和凌晨 1 点在交通模式上非常相似，二者在数值上的距离是 1；但凌晨 23 点和凌晨 0 点在交通模式上同样非常相似，二者在数值上的距离却是 23。线性数值编码把这两对时刻的“相似性”描述得完全相反。

解决方案是使用正弦和余弦编码，将小时数 $h$ 映射为：

\[
  h_{\sin} = \sin\!\left(\frac{2\pi h}{24}\right), \quad h_{\cos} = \cos\!\left(\frac{2\pi h}{24}\right)
\]

这样，23 点和 0 点在这两个维度上的欧氏距离接近 0，而 0 点和 12 点的距离最大，与实际的交通规律一致。同样的逻辑适用于星期（周期为 7）和月份（周期为 12）。

在 Prompt 中说明这一需求时：

\begin{quote}
请对 \texttt{timestamp} 字段提取以下时间特征：

\begin{itemize}
  \item 小时（0--23），并添加 \texttt{hour\_sin} 和 \texttt{hour\_cos} 两列，使用周期为 24 的正弦余弦编码；
  \item 星期几（0 代表周一，6 代表周日），并添加 \texttt{dow\_sin} 和 \texttt{dow\_cos} 两列，使用周期为 7 的正弦余弦编码；
  \item 是否为法定节假日（布尔型），参考[附上节假日列表或说明如何获取]；
  \item 请不要直接将小时数和星期数作为数值特征，只使用编码后的版本。
\end{itemize}
\end{quote}

\subsection{类别特征的编码选择}

类别特征的编码方式选择同样是需要人做判断的决策，不同的编码对不同的模型有不同的效果。

\textbf{独热编码（One-Hot Encoding）}将每个类别展开为一个二值列，适用于无序类别特征（例如检测器类型：感应线圈、微波、视频）且类别数量较少（一般不超过 20）的情形。需要注意，类别数量多时会产生稀疏的高维特征，不适合树模型以外的大多数算法。

\textbf{标签编码（Label Encoding）}将每个类别映射为一个整数，适用于有序类别（例如道路等级：快速路 $>$ 主干道 $>$ 次干道 $>$ 支路），或仅在树模型中使用的情形（树模型可以正确处理整数标签）。对无序类别使用标签编码会让模型误以为类别之间有大小关系，不应在此情形下使用。

\textbf{目标编码（Target Encoding）}将每个类别替换为该类别对应的目标变量均值（例如，将“检测器所在区域”替换为该区域的平均速度），适用于高基数类别（类别数量大）且类别与目标变量有关联的情形。必须只用训练集的目标均值计算编码，否则会引入数据泄露；对低频类别容易产生过拟合，需要使用平滑处理。

在向 AI 描述类别编码需求时，应明确说明编码方式的选择和理由，不要让 AI 自行决定：

\begin{quote}
对以下类别字段进行编码，请按照我指定的方式处理，不要自行调整：

\begin{itemize}
  \item \texttt{road\_type}（快速路/主干道/次干道/支路）：使用标签编码，映射为 3/2/1/0，体现等级顺序；
  \item \texttt{detector\_type}（感应线圈/微波/视频）：使用独热编码，展开为三列，不删除其中一列（不使用 drop\_first）；
  \item \texttt{district\_id}（行政区编号，共 13 个取值）：使用目标编码，目标变量为 \texttt{avg\_speed}，编码参数只从训练集计算，测试集使用训练集的编码映射表。
\end{itemize}

请将编码后的结果存入新列（原列名加 \texttt{\_enc} 或展开后列名），保留原始类别列，并输出编码映射表供核查。
\end{quote}

\subsection{为什么这些决策不能完全交给AI}

标准化和编码的方法选择，表面上是技术问题，本质上是建模决策。它们的选择依赖以下判断：下游模型是什么（树模型对特征缩放不敏感，神经网络敏感）；特征是否有内在的顺序或周期性（决定标签编码 vs 独热编码 vs 周期编码）；数据划分是否已经确定（决定 scaler 的拟合集）；类别频率分布是否均衡（影响独热编码和目标编码的效果）。

这些判断来自对数据和模型的理解，不是自动化可以替代的。AI 可以帮你快速实现你选定的编码方案，生成验证代码，并在编码结果中发现格式错误或边界条件遗漏——但选择本身必须由你完成。

% ----------------------------------------------------------------
\section{实战练习}
% ----------------------------------------------------------------

\subsection*{练习一：判断缺失机制并设计处理规则}

以下是三个交通数据缺失的场景描述，请分别判断其缺失机制（MCAR、MAR 或 MNAR），并说明判断理由，然后设计合理的处理方案：

\begin{enumerate}
  \item 某城市地铁刷卡数据中，部分时段某站点的进站人数为空。进一步调查发现，该站点在下班高峰期（17:00--19:00）的数据丢失率明显高于其他时段，设备维护记录显示该时段读卡机处于超负荷状态。
  \item 出租车 GPS 轨迹数据中，约有 2\% 的记录经纬度字段为空，且这些记录均匀分布于一天各时段和各区域，无明显规律。经查，空值来源于 GPS 信号偶发性丢失。
  \item 道路事故记录中，“受伤人数”字段在约 15\% 的记录中为空。进一步分析发现，空值记录集中于轻微事故类别（财产损失事故），而伤亡事故记录的完整率接近 100\%。
\end{enumerate}

对每种情形，请说明：如果直接删除缺失记录，会对后续分析产生什么影响？你推荐的处理方案是什么，理由是什么？

\subsection*{练习二：设计异常值检测的分层规则}

选取以下任意一种交通数据场景，按本章介绍的分层检测策略，完整设计异常值检测规则：

\begin{enumerate}
  \item 共享单车骑行记录：每条记录包含借车时间、还车时间、起点桩号、终点桩号、骑行距离和骑行时长。
  \item 高速公路收费数据：每条记录包含车牌号、入口站、出口站、入站时间、出站时间和通行费金额。
  \item 城市公交到站时刻记录：每条记录包含线路号、车辆编号、站点编号、计划到站时间和实际到站时间。
\end{enumerate}

你的回答应至少包含：领域规则层的具体阈值和来源（为什么这个值一定是错误的？）；统计方法层的方法选择和参数（为什么选择 IQR 而非 Z-score，或反之？）；上下文验证层的具体检查内容；以及如何区分数据错误和真实极端事件。

然后，为你设计的规则写一段完整的 Prompt，要求 AI 生成对应的检测代码。

\subsection*{练习三：修复一条有问题的清洗Prompt}

以下是一条关于交通流量数据清洗的 Prompt，其中存在几处会导致问题的设计：

\begin{quote}
我有一份路段流量数据，字段包括时间、路段编号、速度和流量。请帮我清洗这份数据：处理好缺失值，删除异常值，统一时间格式，最后告诉我数据质量怎么样。
\end{quote}

请分析这条 Prompt 中存在哪些问题（至少找出 4 处），并将其改写为符合本章要求的高质量版本。改写后的 Prompt 应覆盖：数据说明、处理规则（包括领域知识的说明）、执行约束、计划确认要求和验证摘要要求。

\subsection*{练习四：清洗流程的完整实战}

选取一份真实的交通数据集（例如 PeMS 检测器数据、某城市公交 GPS 数据或共享单车订单数据），按照本章介绍的流程完成完整的数据清洗工作：

\begin{enumerate}
  \item 进行初始审查，记录数据的基本统计信息，识别主要的质量问题；
  \item 确定每类问题的缺失机制，设计处理规则，写成文字说明；
  \item 将清洗规则转化为 Prompt，分步让 AI 生成代码，先要求 AI 描述执行计划，确认后再执行；
  \item 每步完成后验证结果，记录处理前后的关键统计变化；
  \item 整理清洗决策日志，包含每步应用的规则、处理的记录数量，以及关键的人工判断说明。
\end{enumerate}

最终提交物包括：完整的 Prompt 序列（含你的确认或修正说明）、清洗决策日志、清洗前后的数据对比摘要，以及你在清洗过程中遇到的最意外的数据质量问题说明。





